\banner{Rates, pools and the transport of solutes: the quantitative chain}\label{ch:map3-deep}
\section{Temperature and moisture factors of the biological rates}\label{sec:dm-tf}
Every biological rate of the brown area of the map is written as a rate at reference conditions multiplied by dimensionless factors. The temperature factor is given either by a $Q_{10}$ or by the law of Arrhenius,
\begin{equation}\label{eq:dm-arr}\boxed{f_T=Q_{10}^{(T-T_{\mathrm{ref}})/10}\qquad\text{or}\qquad f_T=\exp\Bigl[\frac{E_a}{R}\Bigl(\frac{1}{T_{\mathrm{ref}}}-\frac1T\Bigr)\Bigr]}\end{equation}
where $E_a$ is the activation energy (J\,mol$^{-1}$), $R=8.314$ J\,mol$^{-1}$\,K$^{-1}$ and the temperatures are in kelvin in the second form. The two forms agree near the reference temperature: an activation energy of 50 to 60 kJ\,mol$^{-1}$ corresponds to a $Q_{10}$ of about 2.

The moisture factor $f_\theta$ is zero in a very dry soil (microbes are inactive below about pF 4.5), rises to one between field capacity and a wetter soil, and falls again near saturation, where oxygen diffuses too slowly (Section~\ref{sec:pg-diff}). Decomposition and nitrification follow this rise-and-fall; denitrification, which needs the absence of oxygen, does the opposite and rises sharply above about 80 per cent of the pores filled with water.

\begin{kfig}{The two rate factors of the biology of the soil. Left: the temperature factor for $Q_{10}=2$ (relative to 10 $^\circ$C). Right: the moisture factor of aerobic processes (decomposition, nitrification; solid) and of denitrification (dashed) against the water-filled pore space.}
\draw[axis] (0,0)--(3.8,0);\draw[axis] (0,0)--(0,3.6);\node[lab,below] at (1.8,-0.1) {$T$ ($^\circ$C)};\node[lab,rotate=90] at (-0.4,1.7) {$f_T$};
\foreach \x/\t in {0/0,1.1/10,2.2/20,3.3/30}{\node[lab,below] at (\x,0) {\t};}
\draw[crv,smooth,domain=0:3.4] plot (\x,{0.8*2^((\x-1.1)/1.1)*0.88});
\begin{scope}[xshift=5.2cm]
\draw[axis] (0,0)--(3.8,0);\draw[axis] (0,0)--(0,3.6);\node[lab,below] at (1.8,-0.1) {water-filled pores (\%)};\node[lab,rotate=90] at (-0.4,1.7) {$f_\theta$};
\foreach \x/\t in {0/0,1.7/50,3.4/100}{\node[lab,below] at (\x,0) {\t};}
\draw[crv] (0,0)--(0.5,0.4)--(1.6,2.9)--(2.2,2.9)--(3.4,1.0);
\draw[gd,line width=0.7pt] (0,0.05)--(2.5,0.2)--(3.0,1.6)--(3.4,3.1);
\node[lab] at (1.4,3.2) {aerobic};\node[lab] at (2.45,2.0) {denitrification};
\end{scope}
\end{kfig}

\exbox{Two ways of correcting for temperature}{A decomposition rate is known at 20 $^\circ$C. Find the factor by which it is lower at 10 $^\circ$C with $Q_{10}=2$, and with the law of Arrhenius and $E_a=55$ kJ\,mol$^{-1}$.}\sol
\givenline{$T_1=20$ $^\circ$C $=293.15$ K ; $T_2=10$ $^\circ$C $=283.15$ K ; $Q_{10}=2$ ; $E_a=55$ kJ\,mol$^{-1}$ ; $R=8.314$ J\,mol$^{-1}$\,K$^{-1}$}
\sh{1. With $Q_{10}$}
\wkt\ (Equation~\ref{eq:dm-arr})
\[f_T=Q_{10}^{(T-T_{\mathrm{ref}})/10}=2^{(10-20)/10}=2^{-1}=0.50\Ans\]
\sh{2. With the law of Arrhenius}
The ratio of the rates at the two temperatures is
\[\frac{k_{20}}{k_{10}}=\exp\Bigl[\frac{E_a}{R}\Bigl(\frac1{T_2}-\frac1{T_1}\Bigr)\Bigr]\ti{i}\]
\[\frac1{283.15}-\frac1{293.15}=0.0035317-0.0034112=1.205\times10^{-4}\ \text{K}^{-1}\]
Substituting in equation \rfi{i},
\[\frac{k_{20}}{k_{10}}=\exp\Bigl(\frac{55\,000}{8.314}\times1.205\times10^{-4}\Bigr)=\exp(0.797)=2.22\]
\[\therefore\quad f_T=\frac1{2.22}=0.45\Ans\]
The two corrections differ by only 10 per cent near the reference temperature.

\section{First-order decay of the organic pools}\label{sec:dm-decay}
Each organic pool of Daisy (added organic matter, soil microbial biomass, soil organic matter) loses carbon at a rate proportional to its size:
\begin{equation}\label{eq:dm-dec}\boxed{\frac{\mathrm dC}{\mathrm dt}=-k\,f_T\,f_\theta\,C\quad\Longrightarrow\quad C(t)=C_0\,e^{-kt},\qquad t_{1/2}=\frac{\ln2}{k}}\end{equation}
The pools differ by orders of magnitude in $k$: fresh residues turn over in weeks, the microbial biomass in months and the stabilised organic matter in centuries. Of the carbon that leaves a pool, a fraction $e$ (the substrate-use efficiency, about 0.4 to 0.6) becomes new microbial biomass and the rest is respired as CO$_2$.

\begin{kfig}{The pools of soil organic matter: added organic matter (AOM) and soil organic matter (SOM) are decomposed by the soil microbial biomass (SMB). A fraction $e$ of the decomposed carbon becomes biomass and $1-e$ is respired; nitrogen follows the carbon according to the C:N ratio of each pool.}
\tikzset{pool/.style={draw=khcyan,fill=khfeat,line width=0.6pt,minimum width=2.2cm,minimum height=0.9cm,font=\sffamily\footnotesize,align=center}}
\node[pool] (aom) at (0,2.6) {AOM\\ residues, manure};\node[pool] (smb) at (4,1.3) {SMB\\ microbes};\node[pool] (som) at (0,0) {SOM\\ humus};
\node[draw=khmag,fill=white,circle,minimum size=1.1cm,font=\sffamily\footnotesize] (co2) at (7.6,1.3) {CO$_2$};
\draw[vec] (aom.east)--(smb.north west);\draw[vec] (som.east)--(smb.south west);\draw[vec] (smb.east)--node[above,lab]{$1-e$}(co2.west);
\draw[vec] (smb.south)--++(0,-1.2)--(som.south east);\node[lab] at (3.3,-0.25) {death, stabilisation};
\node[lab] at (2.3,2.4) {$k_{\mathrm{AOM}}$};\node[lab] at (2.1,0.3) {$k_{\mathrm{SOM}}$};
\end{kfig}

\exbox{Half-life and residue left after a season}{Fresh residue decomposes with $k=0.02$ d$^{-1}$ at reference conditions. Find its half-life and the fraction left after 90 days at reference conditions.}\sol
\givenline{$k=0.02$ d$^{-1}$ ; $t=90$ d}
\wkt\ (Equation~\ref{eq:dm-dec})
\[C(t)=C_0\,e^{-kt}\qquad\text{and}\qquad t_{1/2}=\frac{\ln2}{k}\ti{i}\]
\sh{Half-life}
\[t_{1/2}=\frac{0.693}{0.02}=34.7\ \text{d}\Ans\]
\sh{Fraction left after 90 days}
\[\frac{C}{C_0}=e^{-0.02\times90}=e^{-1.8}=0.165\Ans\]
Thus 83.5 per cent of the residue carbon has left the pool after one season; 90 days is $90/34.7=2.6$ half-lives, and $0.5^{2.6}=0.165$ gives the same answer.

\section{Mineralisation or immobilisation: the critical C:N ratio}\label{sec:dm-cn}
When microbes decompose a substrate with carbon-to-nitrogen ratio $CN_s$ they build biomass of a lower ratio $CN_m$ (about 8). Per unit of decomposed carbon they need $e/CN_m$ of nitrogen and the substrate supplies $1/CN_s$. The net release of mineral nitrogen is therefore
\begin{equation}\label{eq:dm-nmin}\boxed{N_{\mathrm{net}}=C_d\Bigl(\frac1{CN_s}-\frac{e}{CN_m}\Bigr),\qquad CN_{\mathrm{crit}}=\frac{CN_m}{e}}\end{equation}
where $C_d$ is the decomposed carbon. A substrate richer in nitrogen than the critical ratio (about 20 for $e=0.4$ and $CN_m=8$) releases mineral nitrogen; a poorer substrate, such as cereal straw with a ratio of 80, takes mineral nitrogen from the soil (\kw{immobilisation}). This is why incorporating straw can cause a temporary nitrogen deficiency of the next crop.

\exbox{Immobilisation by cereal straw}{Five tonnes of straw dry matter per hectare (45 per cent carbon, $CN_s=80$) are incorporated. The microbes have $e=0.4$ and $CN_m=8$. If all the straw carbon is decomposed, how much mineral nitrogen is immobilised?}\sol
\givenline{$W=5000$ kg DM\,ha$^{-1}$ ; carbon content $=0.45$ ; $CN_s=80$ ; $e=0.4$ ; $CN_m=8$}
\par\indent The decomposed carbon is
\[C_d=5000\times0.45=2250\ \text{kg C\,ha}^{-1}\]
\wkt\ the net release of mineral nitrogen is (Equation~\ref{eq:dm-nmin})
\[N_{\mathrm{net}}=C_d\Bigl(\frac1{CN_s}-\frac{e}{CN_m}\Bigr)\ti{i}\]
\sh{Nitrogen supplied and nitrogen needed}
\[\text{Supplied by the straw}=\frac{C_d}{CN_s}=\frac{2250}{80}=28.1\ \text{kg N\,ha}^{-1}\]
\[\text{Needed by the microbes}=\frac{e\,C_d}{CN_m}=\frac{0.4\times2250}{8}=112.5\ \text{kg N\,ha}^{-1}\]
Substituting in equation \rfi{i},
\[N_{\mathrm{net}}=28.1-112.5=-84.4\ \text{kg N\,ha}^{-1}\Ans\]
The negative sign means \emph{immobilisation}: about 84 kg of mineral nitrogen per hectare are taken from the soil by the microbes, because $CN_s=80$ is far above the critical ratio
\[CN_{\mathrm{crit}}=\frac{CN_m}{e}=\frac8{0.4}=20\]

\section{Transport of solutes: convection, dispersion and sorption}\label{sec:dm-cde}
A solute dissolved in the soil water moves with the water (\kw{convection}) and spreads about the mean position by molecular diffusion and by the variation of the velocities in the pores (\kw{hydrodynamic dispersion}). A sorbed solute is held back. For steady flow these processes are summarised by the convection--dispersion equation with retardation,
\begin{equation}\label{eq:dm-cde}\boxed{R\,\frac{\partial C}{\partial t}=D\,\frac{\partial^2C}{\partial z^2}-v\,\frac{\partial C}{\partial z}-\mu\,R\,C,\qquad R=1+\frac{\rho_b K_d}{\theta},\quad v=\frac q\theta,\quad D=\lambda_dv}\end{equation}
where $C$ is the concentration in the soil water, $v$ the pore-water velocity, $D$ the dispersion coefficient with dispersivity $\lambda_d$ (a few centimetres in the field), $K_d$ the sorption coefficient (L\,kg$^{-1}$), $\rho_b$ the bulk density (kg\,L$^{-1}$) and $\mu$ the first-order degradation rate. The centre of a pulse moves at $v/R$ and its spread grows as $\sqrt{2Dt/R}$. Nitrate is not sorbed ($R=1$) and travels with the water; a pesticide with $K_d$ of a few L\,kg$^{-1}$ travels several times slower and is degraded on the way.

\begin{kfig}{Breakthrough of a solute pulse at three depths: the pulse moves down at the velocity $v/R$ and spreads as $\sqrt{2Dt}$; a sorbed solute (dashed, $R>1$) arrives later and more degraded.}
\draw[axis] (0,0)--(7.0,0);\draw[axis] (0,0)--(0,3.4);\node[lab,below] at (3.5,-0.1) {depth $z$};\node[lab,rotate=90] at (-0.4,1.6) {concentration $C$};
\draw[crv,smooth,domain=0:6.6,samples=80] plot (\x,{3.0*exp(-((\x-1.0)^2)/0.18)});
\draw[crv,smooth,domain=0:6.6,samples=80] plot (\x,{1.6*exp(-((\x-3.0)^2)/0.62)});
\draw[crv,smooth,domain=0:6.6,samples=80] plot (\x,{1.1*exp(-((\x-5.0)^2)/1.3)});
\draw[gd,smooth,domain=0:6.6,samples=80] plot (\x,{1.2*exp(-((\x-1.8)^2)/0.35)});
\node[lab] at (1.0,3.2) {$t_1$};\node[lab] at (3.0,1.85) {$t_2$};\node[lab] at (5.0,1.35) {$t_3$};\node[lab] at (2.35,1.25) {sorbed, $t_2$};
\end{kfig}

\exbox{Travel time and the survival of a pesticide}{A soil has $\theta=0.30$, $\rho_b=1.5$ kg\,L$^{-1}$ and a drainage flux of 300 mm\,yr$^{-1}$. Find the time a pesticide needs to travel 1 m when $K_d=0.5$ L\,kg$^{-1}$ and when $K_d=0.05$ L\,kg$^{-1}$, and the fraction that reaches 1 m if its half-life is 30 days (the attenuation factor $AF=\exp(-0.693\,t/t_{1/2})$).}\sol
\givenline{$\theta=0.30$ ; $\rho_b=1.5$ kg\,L$^{-1}$ ; $q=300$ mm\,yr$^{-1}=0.30$ m\,yr$^{-1}$ ; $z=1$ m ; $t_{1/2}=30$ d ; $K_d=0.5$ and $0.05$ L\,kg$^{-1}$}
\wkt\ the pore-water velocity and the retardation factor are (Equation~\ref{eq:dm-cde})
\[v=\frac q\theta,\qquad R=1+\frac{\rho_bK_d}{\theta}\ti{i}\]
and the solute moves at $v/R$, so its travel time to the depth $z$ is
\[t=\frac{z\,R}{v}\ti{ii}\]
Here $v=0.30/0.30=1.0$ m\,yr$^{-1}$.
\sh{1. Strongly sorbed pesticide ($K_d=0.5$ L\,kg$^{-1}$)}
\[R=1+\frac{1.5\times0.5}{0.30}=1+2.5=3.5\]
\[t=\frac{1\times3.5}{1.0}=3.5\ \text{yr}=1278\ \text{d}\Ans\]
\[AF=\exp\Bigl(-0.693\,\frac{t}{t_{1/2}}\Bigr)=\exp\Bigl(-0.693\times\frac{1278}{30}\Bigr)=\exp(-29.5)=1.5\times10^{-13}\Ans\]
\sh{2. Weakly sorbed pesticide ($K_d=0.05$ L\,kg$^{-1}$)}
\[R=1+\frac{1.5\times0.05}{0.30}=1.25,\qquad t=1.25\ \text{yr}=456\ \text{d}\Ans\]
\[AF=\exp\Bigl(-0.693\times\frac{456}{30}\Bigr)=\exp(-10.5)=2.6\times10^{-5}\Ans\]
\[\therefore\quad\frac{AF(0.05)}{AF(0.5)}=\frac{2.6\times10^{-5}}{1.5\times10^{-13}}\approx2\times10^{8}\]
A tenfold weaker sorption raises the fraction that leaches by eight orders of magnitude: leaching is governed by the ratio of the travel time to the half-life.

\exbox{Spread of a nitrate pulse}{Nitrate ($R=1$) moves with a pore velocity of 100 cm\,yr$^{-1}$ in a soil of dispersivity 5 cm. Find the standard deviation of the pulse when its centre reaches 1 m.}\sol
\givenline{$R=1$ ; $v=100$ cm\,yr$^{-1}$ ; $\lambda_d=5$ cm ; $z=100$ cm}
\par\indent The time for the centre of the pulse to reach 1 m is
\[t=\frac{zR}{v}\ti{i}\]
Substituting the values in equation \rfi{i},
\[t=\frac{100\times1}{100}=1\ \text{yr}\]
\wkt\ the dispersion coefficient is $D=\lambda_dv$ and the spread of the pulse is $\sigma=\sqrt{2Dt/R}$. Thus
\[D=5\times100=500\ \text{cm}^2\,\text{yr}^{-1}\]
\[\sigma=\sqrt{2\times500\times1}=\sqrt{1000}=31.6\ \text{cm}\Ans\]
The pulse is about $2\sigma=63$ cm wide at 1 m, so a single sampling depth sees the nitrate of one fertilisation over several months.

\section{Nitrate leaching and nitrous oxide in numbers}\label{sec:dm-nlos}
The nitrate leached below the root zone is the drainage multiplied by the concentration of the drainage water, $N_L=D\,C_N$. With $D$ in mm and $C_N$ in mg N\,L$^{-1}$, $N_L$ in kg N\,ha$^{-1}$ is $D\,C_N/100$. Concentrations are often given as nitrate (mg NO$_3$\,L$^{-1}$); the conversion is $C_N=C_{\mathrm{NO_3}}\times14/62$. Of the nitrogen that is denitrified, a fraction leaves as nitrous oxide, a greenhouse gas 273 times more potent than CO$_2$ over 100 years.

\exbox{Leaching at the drinking-water limit, and the climate cost of denitrification}{(a) A field drains 250 mm in a winter at the European limit of 50 mg NO$_3$\,L$^{-1}$. How much nitrogen is leached? (b) A wet spring denitrifies 20 kg N\,ha$^{-1}$, a tenth of it as N$_2$O. Find the CO$_2$ equivalent.}\sol
\givenline{$D=250$ mm ; $C_{\mathrm{NO_3}}=50$ mg\,L$^{-1}$ ; denitrified N $=20$ kg\,ha$^{-1}$ ; N$_2$O fraction $=0.1$ ; GWP of N$_2$O $=273$}
\sh{(a) Nitrogen leached}
The concentration as nitrogen is
\[C_N=C_{\mathrm{NO_3}}\times\frac{14}{62}\ti{i}\]
Substituting the values in equation \rfi{i},
\[C_N=50\times0.2258=11.3\ \text{mg N\,L}^{-1}\]
\wkt\ $N_L=D\,C_N/100$ with $D$ in mm and $C_N$ in mg\,L$^{-1}$, so
\[N_L=\frac{250\times11.3}{100}=28.2\ \text{kg N\,ha}^{-1}\Ans\]
\sh{(b) Carbon-dioxide equivalent of the nitrous oxide}
\[\text{N}_2\text{O-N}=0.1\times20=2.0\ \text{kg\,ha}^{-1}\]
\[\text{N}_2\text{O}=2.0\times\frac{44}{28}=3.14\ \text{kg\,ha}^{-1}\tb{$\because$ 28 g N in 44 g N$_2$O}\]
\[\text{CO}_2\text{-eq}=3.14\times273=858\ \text{kg\,ha}^{-1}\Ans\]

